\documentclass[11pt]{article}

\usepackage[margin=1in]{geometry}
\usepackage{amsmath,amssymb}
\usepackage{booktabs}
\usepackage{array}
\usepackage{enumitem}
\usepackage{hyperref}
\usepackage{xcolor}
\hypersetup{colorlinks=true,linkcolor=blue,citecolor=blue,urlcolor=blue}

\newcommand{\relu}{\mathrm{relu}}
\newcommand{\clip}{\mathrm{clip}}
\newcommand{\Wmat}{\mathbf{W}}
\newcommand{\xvec}{\mathbf{x}}
\newcommand{\code}[1]{\texttt{#1}}

\title{Perturbation Testing with Rate and Spiking RNNs:\\
A Tutorial Reproduction with Learnable Time Constants}
\author{Perturbation Testing Tutorial}
\date{\today}

\begin{document}
\maketitle

\begin{abstract}
This note documents a minimal tutorial pipeline for \emph{perturbation testing}
of recurrent network models of cortical spontaneous activity, following
Sourmpis et al.\ (eLife 2026). We fit data-constrained recurrent neural networks
(RNNs) to Neuropixels recordings from a single Allen Institute session and ask
whether the fitted models predict the effect of optogenetic stimulation of
parvalbumin-positive (PV) inhibitory neurons --- a perturbation the model never
saw during training. We summarise the four network architectures (a vanilla rate
RNN, a leaky integrate-and-fire spiking RNN, and their low-rank inter-area
variants), the biological constraints that can be switched on (Dale's law,
E/I balance, spectral-radius control, inter-area Dale's law), the training
objective (a trial-averaged plus a trial-matched loss on an oscillation-feature
embedding), and the evaluation metrics (PSTH correlation, Brain-FID, and
trial-matched $R^2$). We additionally describe a \textbf{membrane time
constant} for the low-rank models, parameterised so it stays positive and bounded,
which can be either learned jointly with the weights or held fixed, and is exposed
as a swept hyperparameter in the experiment driver.
\end{abstract}

\section{Overview and goal}
The aim of perturbation testing is to evaluate a model not on the data
distribution it was trained on, but on its response to a causal intervention.
Here the intervention is optogenetic excitation of PV interneurons. A
spontaneous-activity generator is trained \emph{only} on unperturbed
gray-screen activity; at test time the recorded optogenetic light waveform
$i(t)$ is injected into the model's putative inhibitory units, and the model's
predicted population response is compared against the held-out optogenetic
recordings. A model that has captured the right recurrent structure should
generalise to the perturbation; one that has merely memorised the baseline
statistics should not.

This is a minimal reproduction of Sourmpis et al., eLife 2026
\cite{sourmpis2026}, using data from the Allen Institute Visual Coding
Neuropixels dataset \cite{siegle2021}, session \code{829720705}
(\emph{Pvalb-Cre $\times$ Ai32}, functional-connectivity protocol), which
combines PV optogenetic stimulation with simultaneous Neuropixels recordings.
The trial-matching training loss follows Sourmpis et al., NeurIPS 2023
\cite{sourmpis2023}.

\section{Data}
\label{sec:data}
Spike trains from QC-passing units across six visual areas
(\code{VISp}, \code{VISrl}, \code{VISl}, \code{VISal}, \code{VISpm}, \code{VISam})
are binned at $\mathrm{dt} = 10$\,ms. Neurons are sorted so they are grouped by
area, and every neuron carries an area label and a sign label.

\paragraph{Conditions.} Trials are integer-coded: $c = 0$ for spontaneous
(gray-screen) windows, $c \geq 1$ for drifting gratings, and $c = -1$ for
optogenetic-perturbation trials. Training and the in-distribution test set use a
$70/30$ stratified split per condition; the optogenetic trials form a separate
perturbation set. The generators are trained on the spontaneous block only
($c=0$), binarised to spikes.

\paragraph{Neuron signs.} Each unit is labelled excitatory ($+1$), inhibitory
($-1$), or outlier ($0$) from a fast-spiking waveform classifier. Optionally the
``ideal'' PV labels can instead be derived directly from the perturbation
itself (units whose firing rises during the high-light phase of the strongest
raised-cosine trials; flag \code{-{}-ideal-pv-neurons}). Inhibitory/PV units are
the ones \emph{driven} by the opto light at test time; they are excluded from
scoring (the light clamps them identically in data and model, so scoring them
would inflate the metrics).

\paragraph{Perturbation waveforms.} The Allen functional-connectivity protocol
provides three light waveforms --- single \code{pulse}, \code{fast\_pulses}
($2.5$\,ms pulses at $10$\,Hz over $1$\,s), and \code{raised\_cosine} --- each at
three intensity levels $\{1.3, 1.7, 2.0\}$. In-block \code{sham} segments
(LED-free windows drawn from the gaps between light events) act as a matched
control. Each perturbation trial is $1.5$\,s long with light onset at $0.5$\,s.

\section{Models}
\label{sec:models}
All four models are \emph{spontaneous-activity generators}: there is no external
stimulus input; recurrent dynamics are driven only by per-neuron Gaussian noise.
We write the recurrence as $\mathbf{u} = \xvec\Wmat$, so $W_{ij}$ is the weight
from presynaptic $i$ to postsynaptic $j$. All generators warm-start each call
from the previous call's detached end state (no backpropagation across calls).

\subsection{Vanilla rate RNN (\code{rnn})}
A ReLU-like rate unit clipped to $[0,1]$:
\begin{align}
\mathbf{u}_t &= \xvec_{t-1}\Wmat + \mathbf{b} + \boldsymbol{\sigma}\odot\boldsymbol{\varepsilon}_t + \mathbf{p}_t, \\
\xvec_t &= \clip(\mathbf{u}_t,\, 0,\, 1),
\end{align}
where $\boldsymbol{\varepsilon}_t \sim \mathcal{N}(0,I)$, $\boldsymbol{\sigma}$ is
a learnable per-neuron noise scale (initialised to $0.1$ so the noise-driven
network does not collapse to the zero fixed point), and $\mathbf{p}_t$ is the
optional perturbation current (zero during training). The output is scaled by a
learnable scalar.

\subsection{Leaky integrate-and-fire spiking RNN (\code{lif})}
A spiking variant with a leaky membrane and per-synapse transmission delays.
Each synapse $i\!\to\!j$ is assigned at initialisation to one of
\code{num\_delays} delay bins via a one-hot tensor $\Wmat^{(d)}$, so it reads its
presynaptic spike from its own delay in the past:
\begin{align}
\mathbf{u}_t &= \textstyle\sum_{k} \big(\mathbf{z}_{t-k}\big)\big(W^{(d)}_{\cdot\cdot k}\odot \Wmat\big) + \mathbf{b} + \boldsymbol{\sigma}\odot\boldsymbol{\varepsilon}_t + \mathbf{p}_t, \\
\mathbf{v}_t &= \alpha\,\mathbf{v}_{t-1} + (1-\alpha)\,\mathbf{u}_t - \mathbf{z}_{t-1}\,v_{\mathrm{thr}}, \\
p_t &= \sigma_{\mathrm{logistic}}\!\big(\beta(\mathbf{v}_t - v_{\mathrm{thr}})\big), \qquad
\mathbf{z}_t = \mathbb{1}[p_t > 0.5] + (p_t - \mathrm{sg}(p_t)),
\end{align}
where $\alpha$ is the membrane leak/retention factor, $v_{\mathrm{thr}}$ the
firing threshold, and the last term is a straight-through estimator (forward:
sampled $0/1$ spike; backward: gradient through $p_t = \sigma_{\mathrm{logistic}}$).
$\mathrm{sg}(\cdot)$ denotes stop-gradient.

\subsection{Low-rank inter-area rate RNN (\code{lowRNN})}
\label{sec:lowrnn}
This model imposes \emph{area structure} on the connectivity. The global weight
matrix is assembled block-wise: each intra-area block is a full $(N_a \times N_a)$
matrix, while each inter-area block is constrained to rank $R$:
\begin{equation}
\Wmat^{(\mathrm{tgt},\mathrm{src})} =
\begin{cases}
\Wmat_{\text{full}} & \text{intra-area } (\mathrm{tgt}=\mathrm{src}),\\[2pt]
M N^{\top} & \text{inter-area, unconstrained},\\[2pt]
\mathbf{s}_{\mathrm{src}}\odot\big(M N^{\top}\big) & \text{inter-area, Dale-constrained},
\end{cases}
\qquad M \in \mathbb{R}^{N_{\mathrm{src}}\times R},\ N \in \mathbb{R}^{N_{\mathrm{tgt}}\times R}.
\end{equation}
The dense matrix has its diagonal zeroed. Crucially, the \code{lowRNN} dynamics
are now \emph{leaky}, with a learnable time constant (Section~\ref{sec:tau}):
\begin{equation}
\xvec_t = \alpha\,\xvec_{t-1} + (1-\alpha)\,\clip(\mathbf{u}_t, 0, 1),
\qquad
\mathbf{u}_t = \xvec_{t-1}\Wmat_{\text{global}} + \mathbf{b} + \boldsymbol{\sigma}\odot\boldsymbol{\varepsilon}_t + \mathbf{p}_t.
\label{eq:lowrnn-leaky}
\end{equation}
Setting $\alpha \to 0$ recovers the memoryless map $\xvec_t = \clip(\mathbf{u}_t)$
of the vanilla rate RNN; $\alpha \to 1$ makes the rate change arbitrarily slowly.

\subsection{Low-rank spiking RNN (\code{lowBio})}
The biological counterpart of \code{lowRNN}: it combines the low-rank/area-block
connectivity of Section~\ref{sec:lowrnn} with the delayed-synapse LIF dynamics of
\code{lif}. Its membrane leak $\alpha$ is the same learnable, time-constant--derived
quantity used by \code{lowRNN}, so its membrane time constant is learned and
swept as well.

\section{Biological constraints}
The constraints below are applied at initialisation and re-projected after every
optimiser step (\code{apply\_constraint}).

\begin{description}[leftmargin=1.4em]
\item[Dale's law (sign constraint).] With \code{-{}-sign-constrained}, $\Wmat$ is
initialised sign-constrained and E/I-balanced per column (so the net presynaptic
drive into each postsynaptic neuron is zero at init), and the signs are
re-projected after each step (excitatory rows clamped $\geq 0$, inhibitory rows
$\leq 0$). The construction follows the recipe of \cite{sourmpis2023}.
\item[Self-connection removal.] The diagonal of $\Wmat$ is held at zero.
\item[Spectral-radius control.] Because the recurrence has little leak, an
expansive $\Wmat$ would blow the dynamics up over a trial. After each step the
spectral radius is re-pinned to at most $1.2$ (the low-rank factors $M, N$ are
scaled symmetrically by $\sqrt{\cdot}$ so $MN^\top$ scales correctly).
\item[Inter-area Dale's law.] For the low-rank models, \code{-{}-inter-area-dale}
additionally sign-constrains the inter-area blocks (the factors $M, N$ are kept
non-negative and the fixed row-sign vector $\mathbf{s}_{\mathrm{src}}$ supplies
the sign). Without it, intra-area blocks are sign-constrained but inter-area
projections are free.
\end{description}

\section{Membrane time constant: fixed or learnable}
\label{sec:tau}
The membrane time constant $\tau$ of the low-rank models (\code{lowRNN} and
\code{lowBio}) is parameterised in log-space so it is strictly positive for any
real value, and the per-step retention/leak factor $\alpha$ used in
Eqs.~(\ref{eq:lowrnn-leaky}) is derived from it:
\begin{equation}
\tau = \exp(\theta), \qquad \alpha = \exp(-1/\tau) \in (0,1),
\end{equation}
where $\theta$ is the underlying scalar. Larger $\tau$ gives $\alpha$ closer to
$1$, i.e.\ longer memory and slower dynamics. The initial value $\tau_0$ is set by
the \code{-{}-tau-init} flag (units of time bins). Because $\alpha \in (0,1)$ and
$\tau > 0$ are guaranteed by construction, no extra projection is needed in
\code{apply\_constraint}.

The time constant can be either \textbf{learnable} or \textbf{fixed}:
\begin{itemize}[leftmargin=1.6em]
\item \emph{Learnable} (default): $\theta$ is a trained parameter optimised
jointly with the weights; the network learns its own integration timescale,
initialised at $\tau_0$. The learned $\tau$ is logged each epoch.
\item \emph{Fixed} (\code{-{}-fixed-tau}): $\theta$ is held constant as a
non-trainable buffer, so $\tau \equiv \tau_0$ throughout training. Sweeping
$\tau_0$ in this mode performs a controlled scan over the time constant.
\end{itemize}
In the spiking \code{lowBio} model this same $\alpha$ is the membrane leak, so
making it learnable directly learns (or, when fixed, prescribes) the neuronal
integration timescale.

\section{Training objective}
\label{sec:loss}
Models are trained to match recorded spontaneous activity through two
complementary losses computed on low-pass-filtered rasters (a length-$5$ moving
average), then combined with coefficients $\lambda_{\mathrm{ta}}$ and
$\lambda_{\mathrm{tm}}$:
\begin{equation}
\mathcal{L} = \lambda_{\mathrm{ta}}\,\mathcal{L}_{\mathrm{ta}}
            + \lambda_{\mathrm{tm}}\,\mathcal{L}_{\mathrm{tm}}.
\end{equation}

\paragraph{Trial-averaged loss $\mathcal{L}_{\mathrm{ta}}$.} The squared
difference between the trial-averaged (PSTH) low-passed activity of generated and
recorded batches. This roughly corresponds to the paper's per-neuron term.

\paragraph{Trial-matched loss $\mathcal{L}_{\mathrm{tm}}$.} Single trials carry
variability that a PSTH loss ignores. Following \cite{sourmpis2023}, generated
and recorded trials are embedded with a differentiable feature map
$\phi(\cdot)$, an optimal one-to-one assignment between the two sets of trials is
solved on the \emph{detached} cost matrix (Hungarian algorithm), and the MSE over
matched pairs is minimised (gradient flows through the matched feature entries).

\paragraph{Feature map $\phi$.} Both the differentiable loss (torch) and the
numpy evaluation share a single feature pipeline so they cannot drift apart:
(i) average activity within each area to obtain per-area population channels;
(ii) project each channel onto a bank of log-spaced cosine/sine oscillation
bands and take per-band amplitudes plus a DC term; (iii) z-score against the
training-set statistics. Optimisation uses AdamW (learning rate $10^{-3}$, weight
decay $0.1$); \code{apply\_constraint} runs after each step.

\section{Evaluation metrics}
\label{sec:metrics}
Three metrics compare model-generated rasters to held-out recorded rasters of the
same shape:
\begin{enumerate}[leftmargin=1.6em]
\item \textbf{PSTH Pearson $r$} ($\uparrow$): per-neuron Pearson correlation
between trial-averaged firing rates (model vs.\ data). Trials need not be matched.
\item \textbf{Brain-FID} ($\downarrow$): Fr\'echet distance between Gaussian fits
to the distribution of single-trial feature vectors $\phi(z)$ of model vs.\ data
(a distribution-level score, in the spirit of the image FID).
\item \textbf{Trial-matched $R^2$} ($\uparrow$): $R^2$ between optimally matched
single-trial populations on the same feature embedding $\phi$ used by Brain-FID.
\end{enumerate}
Two reference levels frame the scores: \emph{chance baselines} from time-shuffled
and neuron-shuffled rasters, and an \emph{ideal ceiling} from comparing real
train data against real test data. Brain-FID and trial-matched $R^2$ deliberately
share the feature map $\phi$ so the two numbers are consistent.

\paragraph{Perturbation evaluation.} After training, the recorded light waveform
$i(t)$ for each condition is injected as a perturbation current into the driven
(PV/inhibitory) units, scaled by an \emph{opto intensity} factor that is swept
over \code{-{}-opto-intensities}. Metrics are computed on the non-driven units
only. Conditions span \code{sham} and every (waveform, level) combination present
in the session.

\section{Settings and experiment sweep}
\label{sec:settings}
Table~\ref{tab:settings} lists the command-line settings of \code{train\_rnn.py}.
Table~\ref{tab:sweep} summarises the full sweep run by \code{run\_baselines.sh}.

\begin{table}[t]
\centering
\small
\begin{tabular}{@{}lll@{}}
\toprule
Flag & Meaning & Default \\
\midrule
\code{-{}-model} & \code{rnn} / \code{lif} / \code{lowRNN} / \code{lowBio} & \code{rnn} \\
\code{-{}-sign-constrained} & enable Dale's law (sign-constrained, E/I-balanced $\Wmat$) & off \\
\code{-{}-inter-area-dale} & sign-constrain inter-area blocks (low-rank models) & off \\
\code{-{}-rank} & inter-area low-rank $R$ (low-rank models) & $2$ \\
\code{-{}-tau-init} & \textbf{initial/fixed membrane time constant (bins)} & $2.0$ \\
\code{-{}-fixed-tau} & \textbf{hold $\tau$ fixed at \code{-{}-tau-init} instead of learning it} & off (learnable) \\
\code{-{}-num-delays} & number of synaptic delay bins (spiking models) & $3$ \\
\code{-{}-ideal-pv-neurons} & use perturbation-derived PV labels instead of waveform & off \\
\code{-{}-epochs} & training epochs & $20$ \\
\code{-{}-batch-size} & trials per batch & $64$ \\
\code{-{}-lr} & AdamW learning rate & $10^{-3}$ \\
\code{-{}-coeff\_ta} & weight $\lambda_{\mathrm{ta}}$ of the trial-averaged loss & $1.0$ \\
\code{-{}-coeff\_tm} & weight $\lambda_{\mathrm{tm}}$ of the trial-matched loss & $0.3$ \\
\code{-{}-opto-intensities} & perturbation-current scale factors (swept at eval) & $0.1\ 0.5\ 1.0$ \\
\bottomrule
\end{tabular}
\caption{Command-line settings of the training script. The time constant
(\code{-{}-tau-init}) is the newly added learnable, swept parameter.}
\label{tab:settings}
\end{table}

\begin{table}[t]
\centering
\small
\begin{tabular}{@{}lll@{}}
\toprule
Model family & Swept axes & Runs \\
\midrule
\code{rnn} & $\{$unconstrained, sign-constrained$\}$ & $2$ \\
\code{lif} & $\{$unconstrained, sign-constrained$\}$ & $2$ \\
\code{lowRNN}, \code{lowBio} &
$\{$unconstr., sign-constr., inter-Dale$\}\times\text{rank}\{1,2,3,4\}\times\tau_0\{1,2,4,8\}$ &
$2\times 3\times 4\times 4 = 96$ \\
\midrule
\multicolumn{2}{l}{\textbf{Total}} & $100$ \\
\bottomrule
\end{tabular}
\caption{Sweep performed by \code{run\_baselines.sh}. The low-rank families now
sweep the time-constant initialisation $\tau_0$ in addition to rank and the
constraint configuration. Each run additionally evaluates every opto-intensity in
\code{-{}-opto-intensities}. Results are written to
\code{figures/<model>\_<constraint>[\_rank<R>][\_tau<T>]/}.}
\label{tab:sweep}
\end{table}

\section{Conclusion}
The tutorial provides a compact, end-to-end perturbation-testing pipeline:
data-constrained rate and spiking RNNs, optional biological constraints, a
trial-matching training objective, and distribution-level metrics evaluated
under a held-out optogenetic perturbation. The low-rank, area-structured models
(\code{lowRNN}, \code{lowBio}) now carry a learnable membrane time constant that
is both fit by gradient descent and swept across initialisations, letting the
neuronal integration timescale be studied as a first-class hyperparameter
alongside connectivity rank and Dale's-law constraints.

\begin{thebibliography}{9}
\bibitem{sourmpis2026}
G.\ Sourmpis et al.
\newblock Perturbation testing of data-constrained network models.
\newblock \emph{eLife}, 2026.
\newblock \url{https://elifesciences.org/articles/106827}.

\bibitem{sourmpis2023}
G.\ Sourmpis, C.\ Petersen, W.\ Gerstner, and G.\ Bellec.
\newblock Trial matching: capturing variability with data-constrained spiking
neural networks.
\newblock \emph{Advances in Neural Information Processing Systems (NeurIPS)}, 2023.
\newblock \url{https://papers.neurips.cc/paper_files/paper/2023/hash/ec702dd6e83b2113a43614685a7e2ac6-Abstract-Conference.html}.

\bibitem{siegle2021}
J.\ H.\ Siegle et al.
\newblock Survey of spiking in the mouse visual system reveals functional
hierarchy.
\newblock \emph{Nature}, 592:86--92, 2021.
\newblock \url{https://www.nature.com/articles/s41586-020-03171-x}.
\end{thebibliography}

\end{document}
